\chapter{Як упізнати кішку}
% full version: chapters/12_recognition.tex

\pencilfig{12}{\pencilart{12}}{Велика кішка й маленька --- зовсім різні картинки на датчиках ока, а відповідь одна.}

Кішка на картинці може бути ліворуч чи праворуч, близько чи далеко, на сонці чи в тіні. Щоразу на датчики ока падає зовсім інший візерунок, а ви впізнаєте її за частки секунди. У цьому й полягає складність розпізнавання: те саме треба впізнати в тисячах різних виглядів.

\section*{Конвеєр}

Мозок розв'язує задачу конвеєром із приблизно десятка щаблів, і сигнал проходить його за півтори десятої секунди. На кожному щаблі робляться дві речі. Спочатку --- «і»: нейрон спрацьовує, якщо потрібні деталі стоять у потрібному порядку, скажімо, два відрізки сходяться куточком. Потім --- «або по всіх місцях»: наступний нейрон спрацьовує, якщо куточок є будь-де поблизу, байдуже, де саме. Перша операція робить упізнавання розбірливим, друга --- терпимим до зсуву. Щабель за щаблем деталі складаються в частини, частини --- у предмети, а зсуви й масштаби пробачаються.

\pencilfig{112}{%
\foreach \i/\y/\cx/\cy in {1/1.2/0.15/0.3,2/0.3/0.3/0.1,3/-0.6/0.05/0.05}{
  \draw[penlite] (0,\y-0.3) rectangle (0.6,\y+0.3);
  \draw[pcGray, line width=0.8pt] (\cx,\y-0.3+\cy+0.35) -- (\cx,\y-0.3+\cy) -- (\cx+0.3,\y-0.3+\cy);
  \draw[pen=pcBlue, wash=pcBlue] (1.9,\y) circle (0.2);
  \draw[arr] (0.65,\y) -- (1.65,\y);
  \draw[arr=pcBlue] (2.12,\y) -- (3.62,{0.3+0.12*(2-\i)});}
\draw[pen=pcBlue, wash=pcBlue] (3.9,0.3) circle (0.25);
\draw[arr] (4.2,0.3) -- (5.75,0.3); \node[lbl, anchor=south] at (4.97,0.35) {ще щаблі};
% a small cat
\draw[penlite] (6.3,0.3) circle (0.32);
\draw[penlite] (6.03,0.5) -- (6.07,0.82) -- (6.23,0.6); \draw[penlite] (6.57,0.5) -- (6.53,0.82) -- (6.37,0.6);
\node[lbl, anchor=north] at (0.3,-1.0) {картинка}; \node[lbl, anchor=north, align=center] at (1.9,-1.0) {І: куточок\\у цьому місці};
\node[lbl, anchor=north, align=center] at (3.9,-1.0) {АБО: куточок\\будь-де}; \node[lbl, anchor=north] at (6.3,-1.0) {кішка};
}{Один щабель конвеєра: «і» збирає куточок в одному місці, «або» пробачає, де саме він стоїть.}

Саме цю схему скопіювали згорткові нейромережі, які розпізнають картинки в телефоні. І вони відповідають взаємністю: мережі, навчені впізнавати предмети, непогано передбачають, як відповідають нейрони на верхніх щаблях зору.

\section*{Учитися з відео}

Але хто навчав мозок? Ніхто не показував вам мільйон підписаних картинок «кішка». Правдоподібна відповідь --- відео. Світ плине безперервно, і за частку секунди предмет майже не змінюється, а змінюються його положення й освітлення. Якщо зв'язки міцнішають між тим, що видно зараз, і тим, що було мить тому, мережа сама пов'яже різні вигляди одного предмета: повільне --- це предмет, швидке --- це обставини. Є дослід, де предмети непомітно підміняли під час руху очей, --- і нейрони починали вважати різні предмети одним. Як загальне правило це поки гіпотеза.

\section*{Ілюзії --- відбитки механізмів}

Ілюзії --- не поломки, а сліди того, як машину влаштовано. Якщо око довго дивиться на водоспад, а потім на нерухоме каміння, каміння «повзе» вгору: пару каналів «вниз» і «вгору» розбалансовано втомою одного з них. Знаменита сукня, яку одні бачать біло-золотою, а інші синьо-чорною, ділить людей за тим, яке освітлення їхній мозок мовчки припускає. Намальовані «змії, що обертаються» здаються рухомими, бо яскраві ділянки обробляються трохи швидше за тьмяні, і різниця в часі читається як рух. А куб, який то випинається, то западає, --- це дві групи нейронів, що змагаються одна з одною; зміну результату, за однією з моделей, здебільшого вирішує шум.

Цікаво, що нейромережа, навчена передбачати наступний кадр відео, теж «бачить», як обертаються змії. Отже, деякі ілюзії --- неминуча ціна за вміння передбачати.

\section*{Мозок проти нейромережі}

Схожість не повна. Мозкові вистачає небагатьох прикладів, він учиться майже без підписів і витрачає на все двадцять ватів. Нейромережу можна обдурити непомітною для людини підробкою картинки. Щоправда, і людей подібні підробки трохи збивають, якщо показати їх на мить. Де проходить справжня межа між цими двома машинами --- ще з'ясовують.

\fullversion{12}
